\section{Chapter 3: Random Variables and their Distributions}

\subsection{3.1 Random Variables}

\begin{definition}[Random Variable]
Given an experiment with sample space $S$, a \textbf{random variable} (r.v.) $X$ is a function from the sample space $S$ to the real line $\mathbb{R}$:
\[
X: S \to \mathbb{R}.
\]
\end{definition}

\vspace{0.6em}

\begin{intuition}
A random variable is a deterministic numerical summary of the outcome of a random experiment. The randomness comes entirely from the experiment choosing an outcome $s \in S$ according to a probability function $P$. Once $s$ is realized, the random variable crystallizes into a definite real value $X(s)$.
\end{intuition}

\vspace{0.6em}

\begin{example}[Coin Tosses]
Flip a fair coin twice: $S = \{HH, HT, TH, TT\}$.
\begin{itemize}[itemsep=0.3em]
    \item Let $X$ be the number of Heads: $X(HH) = 2$, $X(HT) = X(TH) = 1$, and $X(TT) = 0$.
    \item Let $Y$ be the number of Tails: $Y = 2 - X$, meaning $Y(s) = 2 - X(s)$ for all $s \in S$.
    \item Let $I$ be the indicator that the first toss is Heads: $I(HH) = I(HT) = 1$, and $I(TH) = I(TT) = 0$.
\end{itemize}
Encoding $H = 1$ and $T = 0$, the sample space is $S = \{(1,1), (1,0), (0,1), (0,0)\}$. The functions can be written as:
\[
X(s_1, s_2) = s_1 + s_2, \quad Y(s_1, s_2) = 2 - s_1 - s_2, \quad I(s_1, s_2) = s_1.
\]
\end{example}

\vspace{0.8em}

\subsection{3.2 Distributions and Probability Mass Functions}

\begin{definition}[Discrete Random Variable and Support]
An r.v.~$X$ is \textbf{discrete} if there exists a finite or countably infinite sequence of values $a_1, a_2, \dots$ such that $P(X = a_j \text{ for some } j) = 1$. The set of values $x$ for which $P(X = x) > 0$ is called the \textbf{support} of $X$.
\end{definition}

\vspace{0.6em}

\begin{definition}[Probability Mass Function]
The \textbf{Probability Mass Function} (PMF) of a discrete r.v.~$X$ is the function $p_X$ given by:
\[
p_X(x) = P(X = x).
\]
\end{definition}

\vspace{0.6em}

\begin{theorem}[Valid PMF Criteria]
A function $p_X$ is a valid PMF for a discrete r.v.~$X$ with support $\{x_1, x_2, \dots\}$ if and only if:
\begin{enumerate}[itemsep=0.3em]
    \item \textbf{Nonnegative:} $p_X(x) > 0$ if $x = x_j$ for some $j$, and $p_X(x) = 0$ otherwise.
    \item \textbf{Sums to 1:} $\displaystyle \sum_{j} p_X(x_j) = 1$.
\end{enumerate}
\end{theorem}

\vspace{0.6em}

\begin{example}[PMFs of Two Coin Tosses]
For the two-toss experiment:
\begin{itemize}[itemsep=0.3em]
    \item $p_X(0) = \dfrac{1}{4}, \quad p_X(1) = \dfrac{1}{2}, \quad p_X(2) = \dfrac{1}{4}$.
    \item $p_Y(0) = \dfrac{1}{4}, \quad p_Y(1) = \dfrac{1}{2}, \quad p_Y(2) = \dfrac{1}{4}$.
    \item $p_I(0) = \dfrac{1}{2}, \quad p_I(1) = \dfrac{1}{2}$.
\end{itemize}
$X$ and $Y$ have identical PMFs ($p_X = p_Y$) even though they are different random variables.
\end{example}

\vspace{0.6em}

\begin{example}[Sum of Two Dice]
Roll two fair six-sided dice, producing independent rolls $X$ and $Y$. Let $T = X + Y$. The sample space has 36 equally likely outcomes. The PMF of $T$ is:
\begin{align*}
P(T=2) &= P(T=12) = \frac{1}{36}, & P(T=5) &= P(T=9) = \frac{4}{36}, \\[0.6em]
P(T=3) &= P(T=11) = \frac{2}{36}, & P(T=6) &= P(T=8) = \frac{5}{36}, \\[0.6em]
P(T=4) &= P(T=10) = \frac{3}{36}, & P(T=7) &= \frac{6}{36}.
\end{align*}
The distribution is symmetric: $P(T = t) = P(T = 14 - t)$.
\end{example}

\vspace{0.6em}

\begin{example}[Children in a Household]
Pick a household in the United States uniformly at random. Let $X$ be the number of children in that household. $X$ is a discrete r.v.~with support $\{0, 1, 2, \dots\}$. The probability $P(X = x)$ is proportional to the total number of U.S.~households with $x$ children.
\end{example}

\vspace{0.8em}

\subsection{3.3 Bernoulli and Binomial Distributions}

\begin{definition}[Bernoulli Distribution]
An r.v.~$X$ has a \textbf{Bernoulli distribution} with parameter $p \in (0,1)$, denoted $X \sim \text{Bern}(p)$, if:
\[
P(X = 1) = p \quad \text{and} \quad P(X = 0) = 1 - p.
\]
Its PMF is $P(X = x) = p^x (1-p)^{1-x}$ for $x \in \{0, 1\}$.
\end{definition}

\vspace{0.6em}

\begin{example}[Canadian Lawyers]
Suppose 1 in 200 Canadian adults is a lawyer. Let $X$ indicate whether a randomly selected Canadian adult is a lawyer. Then $X \sim \text{Bern}(1/200)$:
\[
P(X = 1) = \frac{1}{200}, \qquad P(X = 0) = \frac{199}{200}.
\]
\end{example}

\vspace{0.6em}

\begin{story}[Binomial Distribution]
Perform $n$ independent Bernoulli trials, each with success probability $p$. The total number of successes $X$ follows a \textbf{Binomial distribution} with parameters $n$ and $p$, written $X \sim \text{Bin}(n, p)$.
\end{story}

\vspace{0.6em}

\begin{theorem}[Binomial PMF]
If $X \sim \text{Bin}(n, p)$, then:
\[
P(X = k) = \binom{n}{k} p^k (1-p)^{n-k}, \quad k \in \{0, 1, \dots, n\}.
\]
\end{theorem}
\begin{proof}
Any specific sequence of $k$ successes and $n-k$ failures has probability $p^k(1-p)^{n-k}$ by independence. There are $\binom{n}{k}$ such sequences, each corresponding to a choice of which $k$ trials are successes.
\end{proof}

\vspace{0.6em}

\begin{theorem}
If $X \sim \text{Bin}(n, p)$ and $q = 1 - p$, then $n - X \sim \text{Bin}(n, q)$.
\end{theorem}
\begin{proof}
If $X$ is the number of successes, then $n - X$ is the number of failures. Swapping the roles of success and failure replaces $p$ with $q$. Alternatively, setting $Y = n - X$:
\[
P(Y = k) = P(X = n - k) = \binom{n}{n-k} p^{n-k} q^k = \binom{n}{k} q^k p^{n-k}.
\]
\end{proof}

\vspace{0.6em}

\begin{corollary}
If $X \sim \text{Bin}(n, 1/2)$ with $n$ even, the distribution of $X$ is symmetric about $n/2$:
\[
P(X = n/2 + j) = P(X = n/2 - j) \quad \text{for all } j \ge 0.
\]
\end{corollary}

\vspace{0.8em}

\subsection{3.4 Hypergeometric Distribution}

\begin{story}[Hypergeometric Distribution]
An urn contains $w$ white balls and $b$ black balls. We draw $n$ balls without replacement, such that all $\binom{w+b}{n}$ samples are equally likely. The number of white balls $X$ in the sample follows the \textbf{Hypergeometric distribution}, written $X \sim \text{HGeom}(w, b, n)$.
\end{story}

\vspace{0.6em}

\begin{theorem}[Hypergeometric PMF]
If $X \sim \text{HGeom}(w, b, n)$, its PMF is:
\[
P(X = k) = \frac{\displaystyle\binom{w}{k}\binom{b}{n-k}}{\displaystyle\binom{w+b}{n}},
\]
for integers $k$ satisfying $\max(0, n-b) \le k \le \min(n, w)$.
\end{theorem}

\vspace{0.6em}

\begin{example}[Elk Capture-Recapture]
A forest contains $N$ elk. First, $m$ elk are captured, tagged, and released. Later, a sample of $n$ elk are recaptured at random. The number of tagged elk $X$ in the second sample follows:
\[
X \sim \text{HGeom}(m, N - m, n).
\]
Here, tagged elk are the ``white balls'' ($w = m$), untagged elk are the ``black balls'' ($b = N - m$), and the recaptured group is the sample size ($n$).
\end{example}

\vspace{0.6em}

\begin{example}[Aces in a Poker Hand]
A 5-card hand is dealt from a well-shuffled 52-card deck. The number of aces follows $\text{HGeom}(4, 48, 5)$. The probability of getting exactly three aces is:
\[P(X = 3) = \frac{\displaystyle\binom{4}{3}\binom{48}{2}}{\displaystyle\binom{52}{5}} = \frac{4 \times 1128}{2{,}598{,}960} \approx 0.00174.\]
\end{example}

\vspace{0.6em}

\begin{theorem}
$\text{HGeom}(w, b, n)$ and $\text{HGeom}(n, w+b-n, w)$ are identical distributions.
\end{theorem}
\begin{proof}
Classify balls by two sets of tags: Tag 1 (white vs.\ black) and Tag 2 (sampled vs.\ not sampled). Both models count items having both tags. Algebraically:
\[\frac{\displaystyle\binom{w}{k}\binom{b}{n-k}}{\displaystyle\binom{w+b}{n}} = \frac{w!\,b!\,n!\,(w+b-n)!}{k!\,(w+b)!\,(w-k)!\,(n-k)!\,(b-n+k)!} = \frac{\displaystyle\binom{n}{k}\binom{w+b-n}{w-k}}{\displaystyle\binom{w+b}{w}}.\]
\end{proof}

\vspace{0.6em}

\begin{intuition}[Binomial vs.~Hypergeometric]
The Binomial and Hypergeometric distributions both count successes in $n$ trials. The difference lies in whether the trials are independent:
\begin{itemize}[itemsep=0.3em]
    \item \textbf{Binomial:} Sampling \emph{with} replacement. Trials are independent.
    \item \textbf{Hypergeometric:} Sampling \emph{without} replacement. Trials are dependent because each draw alters the composition of the remaining population.
\end{itemize}
\end{intuition}

\vspace{0.8em}

\subsection{3.5 Discrete Uniform Distribution}

\begin{story}[Discrete Uniform Distribution]
Let $C$ be a finite, non-empty set of numbers. An r.v.~$X$ chosen uniformly at random from $C$ follows the \textbf{Discrete Uniform distribution}, written $X \sim \text{DUnif}(C)$.
\end{story}
Its PMF is $P(X = x) = \dfrac{1}{\vert{}C\vert{}}$ for $x \in C$, and for any subset $A \subseteq C$:
\[P(X \in A) = \frac{\vert{}A\vert{}}{\vert{}C\vert{}}.\]

\vspace{0.6em}

\begin{example}[Random Slips of Paper]
A hat contains 100 slips of paper numbered $1, 2, \dots, 100$. Five slips are drawn.
\begin{itemize}[itemsep=0.4em]
    \item \textbf{With Replacement:}
    \begin{enumerate}[itemsep=0.3em]
        \item Let $Y$ be the number of slips $\ge 80$. Then $P(\text{slip } \ge 80) = 21/100 = 0.21$, so $Y \sim \text{Bin}(5, 0.21)$.
        \item The value of the $j$-th draw is $X_j \sim \text{DUnif}(\{1, 2, \dots, 100\})$.
        \item The probability that 100 is drawn at least once:
        \[1 - P(\text{no 100}) = 1 - \left(\frac{99}{100}\right)^5 \approx 0.0490.\]
    \end{enumerate}
    \item \textbf{Without Replacement:}
    \begin{enumerate}[itemsep=0.3em]
        \item Let $Y$ be the number of slips $\ge 80$. Then $Y \sim \text{HGeom}(21, 79, 5)$.
        \item By symmetry, the $j$-th draw unconditionally satisfies $Y_j \sim \text{DUnif}(\{1, 2, \dots, 100\})$.
        \item The probability that 100 is in the sample:
        \[P(100 \in \text{sample}) = \frac{\displaystyle\binom{1}{1}\binom{99}{4}}{\displaystyle\binom{100}{5}} = \frac{5}{100} = 0.05.\]
    \end{enumerate}
\end{itemize}
\end{example}

\vspace{0.8em}

\subsection{3.6 Cumulative Distribution Functions (CDFs)}

\begin{definition}[Cumulative Distribution Function]
The \textbf{Cumulative Distribution Function} (CDF) of an r.v.~$X$ is:
\[F_X(x) = P(X \le x).\]
\end{definition}

\vspace{0.6em}

\begin{theorem}[Properties of Valid CDFs]
A function $F$ is a valid CDF if and only if:
\begin{enumerate}[itemsep=0.3em]
    \item \textbf{Non-decreasing:} $x_1 \le x_2 \implies F(x_1) \le F(x_2)$.
    \item \textbf{Right-continuous:} $\displaystyle \lim_{t \to x^+} F(t) = F(x)$ for all $x \in \mathbb{R}$.
    \item \textbf{Limits at Infinity:} $\displaystyle \lim_{x \to -\infty} F(x) = 0$ and $\displaystyle \lim_{x \to \infty} F(x) = 1$.
\end{enumerate}
\end{theorem}

\vspace{0.8em}

\subsection{3.7 Functions of Random Variables}

\begin{theorem}[PMF of $g(X)$]
For a discrete r.v.~$X$, the PMF of $Y = g(X)$ is:
\[P(g(X) = y) = \sum_{x : g(x) = y} P(X = x).\]
If $g$ is one-to-one, this simplifies to $P(g(X) = g(x)) = P(X = x)$.
\end{theorem}

\vspace{0.6em}

\begin{example}[1D Random Walk]
A particle starts at 0 and takes $n$ independent steps, each $+1$ or $-1$ with equal probability $1/2$. Let $Y$ be the particle's position after $n$ steps.
Let $X \sim \text{Bin}(n, 1/2)$ be the number of rightward steps. The final position is:
\[Y = X(+1) + (n-X)(-1) = 2X - n.\]
Because $g(x) = 2x - n$ is one-to-one, solving for $x$ gives $x = (n+k)/2$. Thus, for integers $k \in \{-n, -n+2, \dots, n\}$:
\[P(Y = k) = P\left(X = \frac{n+k}{2}\right) = \binom{n}{\frac{n+k}{2}} \left(\frac{1}{2}\right)^n.\]
\end{example}

\vspace{0.6em}

\begin{example}[Distance from Origin]
Continuing the random walk, let $D = \vert{}Y\vert{}$ be the distance from the origin, assuming $n$ is even. For $k = 0$:
\[P(D = 0) = P(Y = 0) = \binom{n}{n/2} \left(\frac{1}{2}\right)^n.\]
For even integers $k \in \{2, 4, \dots, n\}$, since $g(y) = \vert{}y\vert{}$ maps both $k$ and $-k$ to $k$:
\[P(D = k) = P(Y = k) + P(Y = -k) = 2 \binom{n}{\frac{n+k}{2}} \left(\frac{1}{2}\right)^n.\]
\end{example}

\vspace{0.6em}

\begin{example}[Maximum of Two Die Rolls]
Roll two fair six-sided dice, $X$ and $Y$. Let $M = \max(X, Y)$.
Using the CDF:
\[P(M \le k) = P(X \le k, Y \le k) = P(X \le k)P(Y \le k) = \left(\frac{k}{6}\right)^2.\]
The PMF is the difference between consecutive CDF values:
\[P(M = k) = P(M \le k) - P(M \le k-1) = \frac{k^2 - (k-1)^2}{36} = \frac{2k - 1}{36}.\]
Evaluating for $k = 1, \dots, 6$ gives:
\[P(M=1) = \frac{1}{36}, \quad P(M=2) = \frac{3}{36}, \quad P(M=3) = \frac{5}{36},\]
\[P(M=4) = \frac{7}{36}, \quad P(M=5) = \frac{9}{36}, \quad P(M=6) = \frac{11}{36}.\]
\end{example}

\vspace{0.8em}

\subsection{3.8 Independence of Random Variables}

\begin{definition}[Independence of Two R.V.s]
Random variables $X$ and $Y$ are \textbf{independent} if:
\[P(X \le x, Y \le y) = P(X \le x)P(Y \le y) \quad \text{for all } x, y \in \mathbb{R}.\]
For discrete r.v.s, this is equivalent to:
\[P(X = x, Y = y) = P(X = x)P(Y = y) \quad \text{for all } x, y.\]
\end{definition}

\vspace{0.6em}

\begin{theorem}[Sum of Independent Binomials]
If $X \sim \text{Bin}(n, p)$ and $Y \sim \text{Bin}(m, p)$ are independent, then:
\[X + Y \sim \text{Bin}(n+m, p).\]
\end{theorem}
\begin{proof}
Conditioning on $X$ and applying the Law of Total Probability:
\begin{align*}
P(X + Y = k) &= \sum_{j=0}^k P(X + Y = k \mid X = j) P(X = j) \\[0.6em]
&= \sum_{j=0}^k P(Y = k - j) P(X = j) \quad (\text{by independence}) \\[0.6em]
&= \sum_{j=0}^k \binom{m}{k-j} p^{k-j} q^{m-k+j} \binom{n}{j} p^j q^{n-j} \\[0.6em]
&= p^k q^{n+m-k} \sum_{j=0}^k \binom{m}{k-j} \binom{n}{j} \\[0.6em]
&= \binom{n+m}{k} p^k q^{n+m-k} \quad (\text{by Vandermonde's Identity}).
\end{align*}
\end{proof}

\vspace{0.6em}

\begin{example}[Matching Pennies]
Two players flip fair pennies independently. Let $X \in \{-1, 1\}$ and $Y \in \{-1, 1\}$ indicate Heads for players A and B. Let $Z = XY$ indicate whether they match.
$X$ and $Y$ are unconditionally independent, but given $Z = 1$, we know $X = Y$, making them conditionally dependent.
\end{example}

\vspace{0.6em}

\begin{example}[Mystery Opponent]
You play two tennis games against one of two twins. Against Twin 1, the win probability is $1/2$; against Twin 2, it is $3/4$. Let $Z$ indicate playing Twin 1 ($P(Z=1) = 1/2$).
Conditional on $Z=1$, game outcomes $X$ and $Y$ are i.i.d.~$\text{Bern}(1/2)$. Conditional on $Z=0$, they are i.i.d.~$\text{Bern}(3/4)$.
Unconditionally, $X$ and $Y$ are dependent: winning the first game makes it more likely that the opponent is the weaker twin, increasing the probability of winning the second.
\end{example}